% Source: https://artofproblemsolving.com/community/c3h1765109_geometrycounting
% Stable identity: aops-topic-1765109-q01
Put $C$ at the origin with $B=(10,0)$ and $A=(0,10)$. The figure consists of the horizontal lines $y=j$ and the vertical lines $x=i$ ($0\le i,j\le9$, including the legs) together with the hypotenuse $x+y=10$. A triangle needs three sides on three pairwise non-parallel lines, so it uses one horizontal line, one vertical line and the hypotenuse.

\begin{asy}
usepackage("amsmath");
size(5cm);
defaultpen(fontsize(10pt));
pair C=(0,0),A=(0,10),B=(10,0);
for(int i=1;i<10;++i){draw((i,0)--(i,10-i),gray);draw((0,i)--(10-i,i),gray);}
fill((3,2)--(8,2)--(3,7)--cycle,paleblue);
draw((3,2)--(8,2)--(3,7)--cycle,blue+linewidth(1));
draw(A--B--C--cycle,linewidth(0.9));
dot((3,2),blue);
label("$A$",A,N);label("$B$",B,E);label("$C$",C,SW);
label("$(i,j)$",(3,2),SW,blue);
label("\mbox{}",(min(currentpicture,user=true)+max(currentpicture,user=true))/2,black); // keeps the text after the figure black when LaTeX typesets the labels inline
\end{asy}

The lines $x=i$ and $y=j$ meet at $(i,j)$, which lies inside the figure (or on a leg) exactly when $i+j\le9$; the triangle is then the one with vertices $(i,j)$, $(10-j,j)$ and $(i,10-i)$, whose legs both have length $10-i-j$. So every triangle in the figure is an isosceles right triangle, and they correspond to the pairs of integers $i,j\ge0$ with $i+j\le9$. Their number is
\[
10+9+8+\cdots+1=\boxed{55}.
\]
